% Part: sets-functions-relations
% Chapter: functions
% Section: inverses

\documentclass[../../../include/open-logic-section]{subfiles}

\begin{document}

\olfileid[fa-IR]{sfr}{fun}{inv}
\olsection{وارون‌های توابع}

\begin{explain}
توابع را همچون نگاشت‌ها در نظر می‌گیریم. پس پرسشی طبیعی دربارهٔ
توابع این است که آیا می‌توان نگاشت را «وارونه کرد». برای
نمونه، تابع پسین $f(x) = x + 1$ را می‌توان وارونه کرد، به
این معنا که تابع $g(y) = y - 1$ کاری را که $f$ انجام می‌دهد «خنثی می‌کند».

اما باید دقت کنیم. گرچه ضابطهٔ~$g$ تابعی
$\Int \to \Int$ تعریف می‌کند، یک \emph{تابع} $\Nat
\to \Nat$ تعریف نمی‌کند، زیرا $g(0) \notin \Nat$. بنابراین حتی در موارد ساده نیز
کاملاً آشکار نیست که آیا می‌توان تابعی را وارونه کرد؛ این امر ممکن است به
دامنه و هم‌دامنه بستگی داشته باشد.

مفهوم وارون یک تابع این نکته را دقیق‌تر بیان می‌کند.
\end{explain}

\begin{defn}
تابع $g \colon B \to A$ یک \emph{وارون} برای تابع $f
\colon A \to B$ است اگر $f(g(y)) = y$ و $g(f(x)) = x$ برای همهٔ $x \in A$
و $y \in B$ برقرار باشند.
\end{defn}

اگر $f$ وارونی مانند~$g$ داشته باشد، اغلب $f^{-1}$ را به‌جای~$g$ می‌نویسیم.

\begin{explain}
اکنون تعیین می‌کنیم که توابع چه هنگام وارون دارند. نامزدی مناسب
برای وارون $f\colon A \to B$ تابع $g\colon B \to A$ است که «تعریف
می‌شود» به‌وسیلهٔ
\[
g(y) = \text{«آن» $x$ که $f(x) = y$ باشد.}
\]
اما گیومه‌های احتیاطی پیرامون «تعریف می‌شود» (و «آن») نشان می‌دهند که
این یک تعریف نیست. دست‌کم، در کلیت کامل همیشه کار
نخواهد کرد. زیرا برای آنکه این تعریف تابعی را مشخص کند، باید یگانه~$x$ای
وجود داشته باشد که $f(x) =
y$---خروجی~$g$ باید به‌طور یکتا مشخص شود. افزون بر این، باید
برای هر $y \in B$ مشخص شده باشد. اگر $x_1$ و $x_2 \in
A$ چنان باشند که $x_1 \neq x_2$ اما $f(x_1) = f(x_2)$، آنگاه $g(y)$ برای
$y = f(x_1) = f(x_2)$ به‌طور یکتا مشخص نخواهد شد. و اگر هیچ~$x$ای
وجود نداشته باشد که $f(x) = y$، آنگاه $g(y)$ اصلاً مشخص نمی‌شود. به
بیان دیگر، برای آنکه $g$ تعریف شود، $f$ باید هم !!{injective} و هم
!!{surjective} باشد.

گام‌به‌گام پیش برویم. پرسش را به دو بخش تقسیم می‌کنیم: با داشتن
تابع~$f\colon A \to B$، چه هنگام تابعی مانند $g\colon B \to A$
وجود دارد که $g(f(x)) = x$؟ چنین $g$ای کاری را که $f$ انجام می‌دهد «خنثی می‌کند» و
\emph{وارون چپِ}~$f$ نامیده می‌شود. دوم، چه هنگام
تابعی مانند $h\colon B \to A$ وجود دارد که $f(h(y)) = y$؟ چنین $h$ای
\emph{وارون راستِ}~$f$ نامیده می‌شود---$f$ کاری را که $h$ انجام می‌دهد «خنثی می‌کند».
\end{explain}

\begin{prop}
اگر $f\colon A \to B$ !!{injective} باشد، آنگاه \emph{وارون
چپی} مانند~$g\colon B \to A$ برای~$f$ وجود دارد چنان‌که $g(f(x)) = x$ برای همهٔ $x
\in A$.
\end{prop}

\begin{proof}
فرض کنید $f\colon A \to B$ !!{injective} باشد. یک $y \in B$ را در نظر بگیرید.
اگر $y \in \ran{f}$، یک $x \in A$ وجود دارد چنان‌که $f(x) = y$. چون
$f$ !!{injective} است، فقط یک چنین~$x \in A$ای وجود دارد. پس می‌توانیم
تعریف کنیم: $g(y) = x$؛ یعنی $g(y)$ «آن» $x \in A$ای است که $f(x)
= y$. اگر $y \notin \ran{f}$، می‌توانیم آن را به هر~$a \in A$ای نگاشت کنیم. پس
می‌توانیم یک $a \in A$ برگزینیم و $g\colon B \to A$ را به‌صورت زیر تعریف کنیم:
\[
g(y) = \begin{cases}
    x & \text{اگر $f(x) = y$}\\
    a & \text{اگر $y \notin \ran{f}$ باشد.}
\end{cases}
\]
این تابع برای همهٔ $y \in B$ تعریف شده است، زیرا برای هر چنین $y \in \ran{f}$ای
دقیقاً یک $x \in A$ وجود دارد چنان‌که $f(x) = y$. بنا بر تعریف، اگر
$y = f(x)$، آنگاه $g(y) = x$؛ یعنی $g(f(x)) = x$.
\end{proof}

\begin{prob}
نشان دهید اگر $f\colon A \to B$ وارون چپی مانند~$g$ داشته باشد، آنگاه $f$
!!{injective} است.
\end{prob}

\begin{prop}
    اگر $f\colon A \to B$ !!{surjective} باشد، آنگاه
    \emph{وارون راستی} مانند~$h\colon B \to A$ برای~$f$ وجود دارد چنان‌که $f(h(y)) =
    y$ برای همهٔ~$y \in B$.
\end{prop}

\begin{proof}
فرض کنید $f\colon A \to B$ !!{surjective} باشد. یک $y \in
B$ را در نظر بگیرید. چون $f$ !!{surjective} است، یک $x_y \in A$ وجود دارد که $f(x_y)
= y$. پس می‌توانیم تعریف کنیم: $h(y) = x_y$؛ یعنی برای هر $y \in B$
یک $x \in A$ برمی‌گزینیم چنان‌که $f(x) = y$؛ چون $f$ !!{surjective} است،
همواره دست‌کم یک مورد برای انتخاب وجود دارد.\footnote{چون $f$
!!{surjective} است، برای هر~$y \in B$ مجموعهٔ $\Setabs{x}{f(x) = y}$
ناتهی است. تعریف ما از~$h$ ایجاب می‌کند که از هریک از این مجموعه‌ها یک $x$
برگزینیم. اینکه این کار همواره ممکن است در واقع
آشکار نیست---امکان انجام این انتخاب‌ها صرفاً به‌عنوان یک اصل موضوع پذیرفته می‌شود.
به بیان دیگر، این گزاره «اصل انتخاب» را مفروض می‌گیرد؛ موضوعی که
\oliflabeldef{sth:choice::chap}{در \olref[sth][choice][]{chap} دوباره بررسی خواهیم کرد}{در اینجا از آن چشم می‌پوشیم}.
بااین‌حال، در بسیاری از موارد خاص، مثلاً هنگامی که $A = \Nat$ یا متناهی است، یا هنگامی که $f$ !!{bijective} است،
به اصل انتخاب نیازی نیست. (در حالت خاصی که $f$ !!{bijective} است، برای هر $y \in B$ مجموعهٔ
$\Setabs{x}{f(x) = y}$ دقیقاً یک !!{element} دارد؛ پس هیچ انتخابی در کار نیست.)}
بنا بر تعریف، اگر $x = h(y)$،
آنگاه $f(x) = y$؛ یعنی برای هر $y \in B$، $f(h(y)) = y$.
\end{proof}

\begin{prob}
نشان دهید اگر $f\colon A \to B$ وارون راستی مانند~$h$ داشته باشد، آنگاه $f$
!!{surjective} است.
\end{prob}

\begin{explain}
  با ترکیب ایده‌های برهان پیشین، اکنون نتیجه می‌گیریم که هر
  !!{bijection} وارونی دارد؛ یعنی تابع واحدی وجود دارد
  که هم وارون چپ و هم وارون راستِ~$f$ است.
\end{explain}

\begin{prop}\ollabel{prop:bijection-inverse}
اگر $f\colon A \to B$ !!{bijective} باشد، تابعی مانند
~$f^{-1}\colon B \to A$ وجود دارد چنان‌که برای همهٔ $x \in A$،
$f^{-1}(f(x)) = x$ و برای همهٔ $y \in B$، $f(f^{-1}(y)) = y$.
\end{prop}

\begin{proof}
تمرین.
\end{proof}

\begin{prob}
گزارهٔ \olref[sfr][fun][inv]{prop:bijection-inverse} را اثبات کنید. باید
~$f^{-1}$ را تعریف کنید، نشان دهید تابع است، و نشان دهید
وارون~$f$ است؛ یعنی $f^{-1}(f(x)) = x$ و $f(f^{-1}(y)) = y$ برای
همهٔ $x \in A$ و $y \in B$.
\end{prob}

\begin{explain}
راهی اندکی کلی‌تر برای استخراج وارون‌ها وجود دارد. در
\olref[kin]{sec} دیدیم که هر تابع $f$، !!a{surjection} $f'
\colon A \to \ran{f}$ را با قرار دادن $f'(x) = f(x)$ برای همهٔ $x \in A$ القا می‌کند.
آشکار است که اگر $f$ !!{injective} باشد، آنگاه $f'$ !!{bijective} است، پس
بنا بر \olref{prop:bijection-inverse} وارون یکتایی دارد. با تسامحی بسیار
جزئی در نمادگذاری، گاهی وارون $f'$ را صرفاً
«وارون~$f$» می‌نامیم.
\end{explain}

\begin{prop}\ollabel{prop:left-right}%
  نشان دهید اگر $f\colon A \to B$ وارون چپی مانند~$g$ و وارون
  راستی مانند~$h$ داشته باشد، آنگاه $h = g$.
\end{prop}

\begin{proof}
  تمرین.
\end{proof}

\begin{prob}
  گزارهٔ \olref[sfr][fun][inv]{prop:left-right} را اثبات کنید.
\end{prob}

\begin{prop}\ollabel{prop:inverse-unique}
هر تابع~$f$ حداکثر یک وارون دارد.
\end{prop}

\begin{proof}
  فرض کنید $g$ و $h$ هر دو وارون~$f$ باشند. آنگاه به‌ویژه
  $g$ وارون چپِ~$f$ و $h$ وارون راست آن است. بنا بر
  \olref{prop:left-right}، $g = h$.
\end{proof}

\end{document}
